\section{Results: single-series corrections}
\label{sec:results}

\begin{table}[t]
\centering
\small
\caption{Size of the single-series corrections: percentage of skill-less searches (20 strategies, 504 daily returns) whose best strategy is called skilled at the nominal 5\% level. Null Zoo v1, 10{,}000 searches per cell.}
\label{tab:size-single}
\begin{tabular}{lrrrrrrrr}
\toprule
Family & LET & LET-Lo & LET-NN & HL & DSR & Boot & HAC-$t$ & SB-$t$ \\
\midrule
IID normal & 4.9 & 5.1 & 4.9 & 4.8 & \textit{0.1} & 4.8 & \textbf{5.6} & \textbf{6.1} \\
Student $t_4$ & 4.8 & 5.2 & \textbf{6.7} & 4.6 & \textit{0.1} & \textbf{6.3} & \textbf{5.7} & \textbf{7.9} \\
Skew $-1.32$ & \textbf{8.3} & \textbf{8.7} & \textit{4.4} & \textbf{8.1} & \textit{0.1} & 5.5 & \textbf{9.0} & \textbf{6.7} \\
Skew $+1.32$ & \textit{3.0} & \textit{3.4} & \textbf{6.1} & \textit{2.9} & \textit{0.1} & 4.9 & \textit{3.7} & \textbf{6.6} \\
GARCH(1,1) & 4.9 & 5.0 & 5.1 & 4.8 & \textit{0.1} & 4.9 & 5.3 & \textbf{6.0} \\
AR(1), $\phi=0.2$ & \textbf{19.5} & 5.3 & \textbf{19.4} & \textbf{19.1} & \textit{0.3} & \textbf{18.4} & \textbf{7.4} & \textbf{7.6} \\
Two vol.\ regimes & 5.0 & 5.2 & 5.5 & 4.9 & \textit{0.1} & 5.4 & 5.1 & \textbf{6.0} \\
Correlated, $\rho=0.5$ & \textit{3.5} & \textit{3.8} & \textit{3.6} & \textit{3.5} & \textit{1.7} & \textit{3.3} & \textit{4.0} & \textit{4.3} \\
Clusters, $\rho=0.8$ & \textit{3.3} & \textit{3.5} & \textit{3.3} & \textit{3.2} & \textit{0.0} & \textit{3.2} & \textit{3.7} & \textit{3.9} \\
\bottomrule
\end{tabular}
\par\smallskip\begin{minipage}{0.97\linewidth}\footnotesize Bold: above the 99\% Monte Carlo band around 5\% (4.44--5.56 for 10{,}000 searches); italic: below it. HL is the Harvey--Liu haircut read one-sided (two-sided $p$ halved); Boot uses 1{,}999 resamples and SB-$t$ 4{,}999.\end{minipage}
\end{table}

\begin{table}[t]
\centering
\small
\caption{Power when one strategy has a true annualized Sharpe ratio of 2, Null Zoo v1 (percent). The oracle knows which strategy is skilled and tests it alone; its power is a natural benchmark, and no best-of-20 test in the table exceeds it.}
\label{tab:power}
\par\textit{Raw power}\par\smallskip
\begin{tabular}{lrrrrrrrr}
\toprule
Family & Oracle & LET & LET-Lo & DSR & Boot & HAC-$t$ & RC & SPA$_c$ \\
\midrule
IID normal & 87.7 & 52.9 & 52.9 & 9.2 & 51.2 & 54.0 & 53.5 & 56.9 \\
Student $t_4$ & 88.7 & 54.0 & 54.0 & 10.3 & 53.3 & 55.5 & 52.0 & 58.2 \\
Skew $-1.32$ & 87.5 & 55.3 & 55.4 & 7.7 & 46.4 & 56.1 & 56.1 & 61.7 \\
Skew $+1.32$ & 89.3 & 52.6 & 53.1 & 11.8 & 58.2 & 53.6 & 51.4 & 54.3 \\
GARCH(1,1) & 87.7 & 54.7 & 54.8 & 11.3 & 53.2 & 55.6 & 51.5 & 58.4 \\
AR(1), $\phi=0.2$ & 77.3 & 60.3 & 35.4 & 8.0 & 58.4 & 40.3 & 38.5 & 42.5 \\
Two vol.\ regimes & 87.9 & 55.8 & 55.5 & 11.7 & 54.5 & 56.7 & 50.4 & 58.6 \\
Correlated, $\rho=0.5$ & 88.1 & 51.8 & 51.6 & 24.7 & 50.4 & 52.3 & 59.2 & 61.0 \\
Clusters, $\rho=0.8$ & 88.4 & 51.9 & 51.9 & 15.5 & 50.2 & 52.9 & 59.9 & 62.1 \\
\bottomrule
\end{tabular}
\par
\medskip
\par\textit{Size-adjusted power}\par\smallskip
\begin{tabular}{lrrrrrrrr}
\toprule
Family & Oracle & LET & LET-Lo & DSR & Boot & HAC-$t$ & RC & SPA$_c$ \\
\midrule
IID normal & 86.7 & 53.2 & 52.7 & 40.4 & 51.2 & 52.4 & 53.2 & 50.2 \\
Student $t_4$ & 88.1 & 54.6 & 53.5 & 41.9 & 45.8 & 53.5 & 52.0 & 51.7 \\
Skew $-1.32$ & 84.9 & 47.1 & 46.4 & 39.2 & 42.8 & 46.0 & 54.8 & 44.6 \\
Skew $+1.32$ & 90.0 & 60.0 & 58.6 & 43.0 & 58.2 & 57.6 & 51.7 & 56.4 \\
GARCH(1,1) & 87.6 & 54.9 & 54.8 & 42.9 & 53.2 & 54.6 & 51.8 & 53.8 \\
AR(1), $\phi=0.2$ & 75.1 & 35.4 & 34.2 & 29.8 & n/a & 34.3 & 34.8 & 32.4 \\
Two vol.\ regimes & 87.4 & 55.6 & 54.9 & 43.5 & 51.3 & 56.1 & 50.8 & 56.5 \\
Correlated, $\rho=0.5$ & 87.5 & 56.8 & 55.8 & 40.8 & 54.9 & 56.1 & 56.9 & 54.2 \\
Clusters, $\rho=0.8$ & 88.1 & 58.9 & 58.1 & 51.3 & 57.0 & 58.1 & 58.7 & 55.8 \\
\bottomrule
\end{tabular}
\par
\par\smallskip\begin{minipage}{0.97\linewidth}\footnotesize Size-adjusted power is power at the cutoff at which the same test rejects at most 5\% of the matching null cell; a user who does not know the return family cannot apply it, but it shows how well each statistic ranks skilled searches above skill-less ones. n/a: the test's smallest attainable $p$-value already occurs in more than 5\% of the null cell, so no cutoff attains the level. The Monte Carlo standard error of raw power is at most 0.5 points. Size-adjusted power also carries the error of the estimated cutoff, and the joint tests' $p$-values lie on a grid of $1/1{,}000$, so their achieved null rejection rate at the cutoff can be below 5\%.\end{minipage}
\end{table}



Table~\ref{tab:size-single} reports the size of the single-series corrections in v1 and Table~\ref{tab:power} raw and size-adjusted power at a Sharpe ratio of 2; the appendix gives the remaining validators, the other skill levels and the confirmation run v1b. Rates in the text are from v1 unless stated.

\subsection{The baseline is calibrated, and the search itself costs power}
Under independent normal returns the luck-equivalent-trials test is exact by construction: the $t$-statistic is Student $t$ and the trials are independent. Its size is \nz{v1.size.luck_trials.iid_normal}\% (\nz{v1b.size.luck_trials.iid_normal}\% in v1b). The non-normal-error version (\nz{v1.size.luck_trials_nonnormal_se.iid_normal}\%), the one-sided Harvey--Liu haircut (\nz{v1.size.haircut_bonferroni_one_sided.iid_normal}\%), the bootstrap (\nz{v1.size.bootstrap_best.iid_normal}\%) and the Reality Check (\nz{v1.size.reality_check.iid_normal}\%) are calibrated too. The tests that estimate a long-run variance are slightly liberal even here (HAC-$t$ \nz{v1.size.hac_t_sidak.iid_normal}\%, the single-strategy oracle \nz{v1.size.oracle_hac_t.iid_normal}\%, SB-$t$ \nz{v1.size.bootstrap_t_sidak.iid_normal}\%): a kernel or block estimate of the variance is noisier than the sample variance, a theme Section~\ref{sec:snooping} returns to.

Selection, not the correction, costs most of the power. With a true Sharpe ratio of 2 over two years, the oracle that tests the skilled strategy alone rejects in \nz{v1.power2.oracle_hac_t.iid_normal}\% of searches; the best calibrated best-of-\nz{trials} tests reach \nz{v1.power2.luck_trials.iid_normal}\% (LET) and \nz{v1.power2.reality_check.iid_normal}\% (RC). The skilled strategy is the in-sample best in only \nz{v1.best_skilled2.iid_normal}\% of these searches, so even a perfect correction would crown the wrong strategy in about one search in five; under AR(1) returns it is the best in \nz{v1.best_skilled2.ar1}\%.

\subsection{Autocorrelation breaks the tests that ignore it}
With a lag-one autocorrelation of 0.2, the tests that treat returns as independent rejected about four times as often as their level: \nz{v1.size.luck_trials.ar1}\% for the luck-equivalent-trials test, \nz{v1.size.haircut_bonferroni_one_sided.ar1}\% for the one-sided haircut, \nz{v1.size.luck_trials_nonnormal_se.ar1}\% with the non-normal standard error and \nz{v1.size.bootstrap_best.ar1}\% for the i.i.d.\ bootstrap. Positive autocorrelation makes the sample Sharpe ratio more variable than the i.i.d.\ formula assumes; the non-normal correction adjusts for the shape of the marginal distribution, which is normal here; and the i.i.d.\ bootstrap destroys the very dependence that inflates the sampling variance of the Sharpe ratio. The deflated Sharpe ratio, which also ignores autocorrelation, stays far below its level (\nz{v1.size.deflated_sharpe.ar1}\%) only because it rejects almost nothing anywhere.

Deflating the $t$-statistic by Lo's annualisation factor removes the excess here: the luck-equivalent-trials test falls to \nz{v1.size.luck_trials_lo.ar1}\% (\nz{v1b.size.luck_trials_lo.ar1}\% in v1b). But the AR(1) family is the factor's best case---it satisfies the assumed geometric decay exactly---and in the eight families without autocorrelation the factor raises the size by \nz{v1.lo_shift_range} points, because the estimated factor adds noise to the statistic. The remedies that make no parametric assumption about the dependence reduce the excess without removing it at this sample size: a Newey--West standard error (HAC-$t$, \nz{v1.size.hac_t_sidak.ar1}\%) and a stationary-bootstrap $t$ test (SB-$t$, \nz{v1.size.bootstrap_t_sidak.ar1}\%), and the joint resampling tests, whose blocks carry the dependence, are close but above the band (RC \nz{v1.size.reality_check.ar1}\%, \nz{v1b.size.reality_check.ar1}\% in v1b).

\subsection{Skew and fat tails}
The Student-$t$ tests are liberal under negative skew (\nz{v1.size.luck_trials.skew_negative}\% for the luck-equivalent-trials test, \nz{v1.size.haircut_bonferroni_one_sided.skew_negative}\% for the one-sided haircut, \nz{v1.size.hac_t_sidak.skew_negative}\% for HAC-$t$) and conservative under positive skew (\nz{v1.size.luck_trials.skew_positive}\%). With $T=\nz{obs}$, the $t$-statistic of a skewed series is itself skewed: under negative skew a high sample mean tends to come with a low sample variance, so the upper tail of $t$ is heavier than Student $t$, and the reverse under positive skew.

The non-normal standard error of \citet{mertens2002comments} moves the error rather than removing it: it is conservative under negative skew (\nz{v1.size.luck_trials_nonnormal_se.skew_negative}\%) and liberal under positive skew (\nz{v1.size.luck_trials_nonnormal_se.skew_positive}\%) and fat tails (\nz{v1.size.luck_trials_nonnormal_se.student_t4}\%). It was not designed to correct the null distribution: at a Sharpe ratio of zero its variance term equals one for every i.i.d.\ distribution, and skewness and kurtosis enter only because the term is evaluated at the estimate. Selection then makes that estimate misleading. Under $t_4$ returns, a symmetric distribution, the selected best series has a mean sample skewness of \nz{v1.skew_best.student_t4}: a series tends to win because it had a large positive outlier. The single-strategy bootstrap tests inherit the same selection effect (Boot \nz{v1.size.bootstrap_best.student_t4}\%, SB-$t$ \nz{v1.size.bootstrap_t_sidak.student_t4}\% under $t_4$): resampling a winner that owes its rank to an outlier reproduces the outlier's influence. No single-series validator is calibrated across negative skew, positive skew and fat tails; the Reality Check, which resamples all strategies jointly, is (\nz{v1.size.reality_check.skew_negative}\%, \nz{v1.size.reality_check.skew_positive}\% and \nz{v1.size.reality_check.student_t4}\%).

Under GARCH volatility clustering and regime switching the Student-$t$ tests stay near 5\% (LET \nz{v1.size.luck_trials.garch}\% and \nz{v1.size.luck_trials.regimes}\%): returns remain a martingale difference sequence with finite variance, and two years of daily data are enough for the $t$-test to be nearly exact.

\subsection{Correlated strategies make the Bonferroni-type tests conservative}
When every pair of strategies has correlation 0.5, or strategies come in clusters of five with correlation 0.8, the search contains fewer than \nz{trials} effectively independent trials, and every test that assumes \nz{trials} is conservative: \nz{v1.size.luck_trials.correlated_trials}\% and \nz{v1.size.luck_trials.block_cluster}\% for the luck-equivalent-trials test, \nz{v1.size.bootstrap_best.correlated_trials}\% and \nz{v1.size.bootstrap_best.block_cluster}\% for the bootstrap. The unused level costs power (LET \nz{v1.power2.luck_trials.correlated_trials}\% raw against \nz{v1.adj2.luck_trials.correlated_trials}\% size-adjusted in the correlated family). The Reality Check, which resamples the strategies together and so sees their correlation, stays close to its level (\nz{v1.size.reality_check.correlated_trials}\% and \nz{v1.size.reality_check.block_cluster}\%; \nz{v1b.size.reality_check.correlated_trials}\% and \nz{v1b.size.reality_check.block_cluster}\% in v1b).

\subsection{The deflated Sharpe ratio at 0.95 tests a harder hypothesis}
\citet{bailey2014dsr} accept a strategy when $\mathrm{DSR}\ge 0.95$. Scored this way, the deflated Sharpe ratio rejects at most \nz{v1.size.deflated_sharpe.correlated_trials}\% of skill-less searches (in the correlated family) and fails to accept \nz{v1.dsr_fail_to_accept.iid_normal}\% of i.i.d.\ normal searches that contain a strategy with a true Sharpe ratio of 2 over two years (\nz{v1.dsr_fail_to_accept_range}\% across families; power \nz{v1.power2.deflated_sharpe.iid_normal}\% against \nz{v1.power2.luck_trials.iid_normal}\% for the \v{S}id\'ak test).

The hurdle behind this conservatism, and how much of the shortfall comes from the threshold rather than the statistic, are worked out in Appendix~\ref{app:dsr}.

Two further results are in Appendix~\ref{app:more}: a bootstrap test's power depends on its number of resamples (with 400 resamples and a \v{S}id\'ak adjustment, power under i.i.d.\ normal returns was \nz{br.power.b400.iid_normal}\%, against \nz{br.power.b1999.iid_normal}\% with 1,999), and the Harvey--Liu haircut read one-sided is the \v{S}id\'ak test up to the Bonferroni approximation.
